% Problema: límites de cocientes (castellano).
%
% Enunciado, pista, resultado, solución y corrección en UN fichero. Cada nivel
% sale solo en el perfil que le toca: la hoja del alumno lleva el enunciado y
% la pista; la hoja con resultados, además, el resultado; la copia del
% profesor, todo. No hay un fichero de soluciones aparte que se pueda quedar
% desincronizado.

\begin{exercise}[Límites de cocientes]
Calcula los siguientes límites:
\[
  \text{(i)}\ \lim_{x \to 2} \frac{x^2 - 4}{x - 2} ,
  \qquad
  \text{(ii)}\ \lim_{x \to 1} \frac{x^3 - 1}{x^2 - 1} ,
  \qquad
  \text{(iii)}\ \lim_{x \to 0} \frac{\sqrt{1 + x} - 1}{x} .
\]

\dmarks{3}

\begin{hint}
En los tres, el numerador y el denominador se anulan en el punto. Factoriza
para simplificar el factor común; en el último, multiplica y divide por
$\sqrt{1 + x} + 1$.
\end{hint}

\begin{answer}
(i) $4$; (ii) $3/2$; (iii) $1/2$.
\end{answer}

\begin{solution}
Como el límite no depende del valor en el punto, se puede simplificar para
$x$ distinto de él y evaluar después.
\begin{parts}
  \item Para $x \neq 2$, $\dfrac{x^2 - 4}{x - 2} = x + 2$, así que el límite
        es $4$.
  \item Para $x \neq \pm 1$,
        \[
          \frac{x^3 - 1}{x^2 - 1}
          = \frac{(x - 1)(x^2 + x + 1)}{(x - 1)(x + 1)}
          = \frac{x^2 + x + 1}{x + 1} ,
        \]
        que tiende a $3/2$.
  \item Para $x \neq 0$,
        \[
          \frac{\sqrt{1 + x} - 1}{x}
          = \frac{(1 + x) - 1}{x\,\bigl(\sqrt{1 + x} + 1\bigr)}
          = \frac{1}{\sqrt{1 + x} + 1} ,
        \]
        que tiende a $1/2$.
\end{parts}
\end{solution}

\begin{marking}
Un punto por apartado. Escribir «$0/0$» como resultado no puntúa;
simplificar bien y equivocarse al evaluar, medio punto.
\end{marking}
\end{exercise}
